\chapter{B1: Nebensätze mit Konjunktionen}
\label{cha:B1Nebensaetze}

Nebensätze werden durch Konjunktionen eingeleitet und haben das Verb am Ende.

\section{Kausalsätze (Grund und Folge)}
\label{sec:B1Kausal}

\begin{center}
\begin{tabular}{|l|c|c|}
\hline
Konjunktion & Position & Beispiel \\ \hline
weil & Verb am Ende & Ich lerne, \textbf{weil} ich lernen muss. \\ \hline
da & Verb am Ende & \textbf{Da} ich müde war, ging ich ins Bett. \\ \hline
denn & Hauptsatz! & Ich lerne, \textbf{denn} ich will lernen. \\ \hline
\end{tabular}
\end{center}

\begin{infobox}{Unterschied}
\textbf{weil} und \textbf{da} = Nebensatz (Verb am Ende)\\
\textbf{denn} = Hauptsatz (Verb an Position 2)
\end{infobox}

\subsection*{Übung 1: weil, denn oder da?}
\begin{enumerate}
    \item Ich bleibe zu Hause, \underline{\hspace{1cm}} es regnet. (weil/da)
    \item \underline{\hspace{1cm}} es schon spät ist, müssen wir gehen. (weil/da)
    \item Er lernt viel, \underline{\hspace{1cm}} er will bestehen. (weil/denn)
    \item \underline{\hspace{1cm}} der Unterricht ausfällt, haben wir frei. (da/weil)
    \item Sie kam zu spät, \underline{\hspace{1cm}} der Zug Verspätung hatte. (weil/da)
\end{enumerate}

\loesung{
\begin{enumerate}
    \item weil
    \item Da
    \item weil/denn (beide möglich)
    \item Da
    \item weil
\end{enumerate}
}

\section{Konzessivsätze (Gegensatz)}
\label{sec:B1Konzessiv}

\begin{center}
\begin{tabular}{|l|c|c|}
\hline
Konjunktion & Position & Beispiel \\ \hline
obwohl & Verb am Ende & \textbf{Obwohl} es regnet, gehen wir spazieren. \\ \hline
trotzdem & Hauptsatz! & Es regnet, \textbf{trotzdem} gehen wir spazieren. \\ \hline
\end{tabular}
\end{center}

\begin{infobox}{Merke}
obwohl = Nebensatz (Verb am Ende)\\
trotzdem = Adverb im Hauptsatz
\end{infobox}

\subsection*{Übung 2: obwohl oder trotzdem?}
\begin{enumerate}
    \item Es war kalt, \underline{\hspace{1cm}} wir sind schwimmen gegangen. (obwohl/trotzdem)
    \item \underline{\hspace{1cm}} er krank war, ist er zur Arbeit gegangen. (obwohl/trotzdem)
    \item Der Film war langweilig, \underline{\hspace{1cm}} sind wir bis zum Ende geblieben. (obwohl/trotzdem)
    \item \underline{\hspace{1cm}} sie müde war, hat sie weitergearbeitet. (obwohl/trotzdem)
    \item Er hatte kein Geld, \underline{\hspace{1cm}} kaufte er ein teures Geschenk. (obwohl/trotzdem)
\end{enumerate}

\loesung{
\begin{enumerate}
    \item obwohl
    \item Obwohl
    \item trotzdem
    \item Obwohl
    \item obwohl (oder trotzdem - unterschiedliche Bedeutung!)
\end{enumerate}
}

\section{Temporalsätze (Zeit)}
\label{sec:B1Temporal}

\begin{center}
\begin{tabular}{|l|c|c|}
\hline
Konjunktion & Bedeutung & Beispiel \\ \hline
bevor & vorher & \textbf{Bevor} ich esse, wasche ich die Hände. \\ \hline
nachdem & nachher & \textbf{Nachdem} ich gegessen habe, gehe ich. \\ \hline
während & gleichzeitig & \textbf{Während} ich arbeite, höre ich Musik. \\ \hline
sobald & sobald & \textbf{Sobald} ich kann, komme ich. \\ \hline
bis & bis & \textbf{Bis} ich fertig bin, warte ich. \\ \hline
\end{tabular}
\end{center}

\begin{tippbox}{Merke bei "nachdem"}
- Vorzeitigkeit: Perfekt/Plusquamperfekt\\
- \textbf{Nachdem} ich \underline{gegessen habe}, gehe ich. (zuerst essen, dann gehen)\\
- \textbf{Bevor} ich \underline{gehe}, esse ich. (erst essen, dann gehen)
\end{tippbox}

\subsection*{Übung 3: temporale Konjunktionen}
\begin{enumerate}
    \item \underline{\hspace{1cm}} ich zur Uni gehe, frühstücke ich. (bevor/nachdem)
    \item \underline{\hspace{1cm}} du kommst, ruf mich an. (sobald/bis)
    \item \underline{\hspace{1cm}} ich las, schlief ich ein. (während/bevor)
    \item \underline{\hspace{1cm}} der Unterricht beginnt, muss ich da sein. (bis/bevor)
    \item \underline{\hspace{1cm}} sie gearbeitet hat, ging sie nach Hause. (nachdem/bevor)
\end{enumerate}

\loesung{
\begin{enumerate}
    \item Bevor
    \item Sobald
    \item Während
    \item Bevor
    \item Nachdem
\end{enumerate}
}

\section{Konditionalsätze (Bedingung)}
\label{sec:B1Konditional}

\begin{center}
\begin{tabular}{|l|c|c|}
\hline
Konjunktion & Bedeutung & Beispiel \\ \hline
wenn & wenn/falls & \textbf{Wenn} es regnet, bleibe ich zu Hause. \\ \hline
falls & falls & \textbf{Falls} du Zeit hast, ruf mich an. \\ \hline
\end{tabular}
\end{center}

\begin{infobox}{Unterschied}
\textbf{wenn} = Realer Bedingung (es regnet wirklich)\\
\textbf{falls} = Hypothetische Bedingung (vielleicht regnet es)
\end{infobox}

\chapterend

\chapter{B1: Passiv}
\label{cha:B1Passiv}

\section{Vorgangspassiv}
\label{sec:B1Vorgangspassiv}

Bildung: \textbf{werden} + Partizip II

\begin{center}
\begin{tabular}{|c|c|c|}
\hline
Aktiv & Passiv & \\ \hline
Der Lehrer erklärt die Regel. & Die Regel wird erklärt. & Präsens \\ \hline
Der Lehrer erklärte die Regel. & Die Regel wurde erklärt. & Präteritum \\ \hline
Der Lehrer hat die Regel erklärt. & Die Regel ist erklärt worden. & Perfekt \\ \hline
\end{tabular}
\end{center}

\warnbox{Achtung}
Bei Perfekt mit Passiv: \textbf{worden} (nicht "geworden")!
\end{warnbox}

\section{Zustandspassiv}
\label{sec:B1Zustandspassiv}

Bildung: \textbf{sein} + Partizip II

\begin{center}
\begin{tabular}{|c|c|c|}
\hline
Vorgangspassiv & Zustandspassiv & \\ \hline
Das Haus wird gebaut. & Das Haus ist gebaut. & (Zustand nach dem Bau) \\ \hline
Die Tür wurde geöffnet. & Die Tür ist geöffnet. & (Zustand nach dem Öffnen) \\ \hline
\end{tabular}
\end{center}

\begin{infobox}{Unterschied}
Vorgangspassiv: betont die Handlung (wer)\\
Zustandspassiv: betont den Zustand (sein)
\end{infobox}

\section{Passiversatz}
\label{sec:B1Passiversatz}

\begin{center}
\begin{tabular}{|l|l|l|}
\hline
Form & Beispiel & Umschreibung \\ \hline
man & Das Haus wird gebaut. & Man baut das Haus. \\ \hline
sein + zu + Infinitiv & Das Haus ist zu bauen. & Das Haus muss gebaut werden. \\ \hline
lassen & Er lässt das Auto reparieren. & Er beauftragt jemanden... \\ \hline
\end{tabular}
\end{center}

\section{Übungen zum Passiv}
\label{sec:B1PassivUebungen}

\subsection*{Übung 1: Aktiv $\rightarrow$ Passiv}
\begin{enumerate}
    \item Der Chef unterschreibt den Vertrag. $\rightarrow$ \underline{\hspace{2cm}}
    \item Die Firma baut ein neues Haus. $\rightarrow$ \underline{\hspace{2cm}}
    \item Die Studenten machen die Übungen. $\rightarrow$ \underline{\hspace{2cm}}
    \item Der Arzt untersucht den Patienten. $\rightarrow$ \underline{\hspace{2cm}}
    \item Die Polizei verhaftet den Dieb. $\rightarrow$ \underline{\hspace{2cm}}
\end{enumerate}

\loesung{
\begin{enumerate}
    \item Der Vertrag wird vom Chef unterschrieben.
    \item Ein neues Haus wird von der Firma gebaut.
    \item Die Übungen werden von den Studenten gemacht.
    \item Der Patient wird vom Arzt untersucht.
    \item Der Dieb wird von der Polizei verhaftet.
\end{enumerate}
}

\chapterend

\chapter{B2: Konjunktiv II - Höflichkeit und Irrealität}
\label{cha:B2KonjunktivII}

\section{Höfliche Bitten und Wünsche}
\label{sec:B2Hoeflich}

\begin{center}
\begin{tabular}{|l|c|c|}
\hline
Normal & Höflich (Konjunktiv II) & \\ \hline
Komm bitte! & Könntest du bitte kommen? & \\ \hline
Hilf mir! & Würdest du mir helfen? & \\ \hline
Ich will wissen... & Ich würde gern wissen... & \\ \hline
Gib mir das Buch! & Würdest du mir das Buch geben? & \\ \hline
\end{tabular}
\end{center}

\begin{infobox}{Merke}
Die常用höflichen Formen mit Konjunktiv II:
- Könnten Sie...?
- Würden Sie...?
- Hätten Sie...?
- Ich würde gern...
\end{infobox}

\section{Ireale Bedingungssätze}
\label{sec:B2Irreal}

\begin{center}
\begin{tabular}{|l|c|c|}
\hline
Bedingung & Folge & \\ \hline
Wenn ich Zeit \underline{hätte}, & \underline{würde} ich kommen. & \\ \hline
Wäre ich reich, & \underline{hätte} ich ein großes Haus. & \\ \hline
Hättest du angerufen, & \underline{wäre} ich zu Hause geblieben. & \\ \hline
\end{tabular}
\end{center}

\begin{tippbox}{Merke}
Irreale Bedingungssätze verwenden Konjunktiv II in beiden Teilen!
\end{tippbox}

\section{Konjunktiv II der wichtigsten Verben}
\label{sec:B2KonjIIVerben}

\begin{center}
\begin{tabular}{|l|c|c|c|}
\hline
Infinitiv & Präteritum & Konjunktiv II & Bedeutung \\ \hline
sein & war & wäre & were \\ \hline
haben & hatte & hätte & had \\ \hline
werden & wurde & würde & would \\ \hline
können & konnte & könnte & could \\ \hline
müssen & musste & müsste & must (subj) \\ \hline
dürfen & durfte & dürfte & may (subj) \\ \hline
wollen & wollte & wollte & wanted (unhöflich!) \\ \hline
mögen & mochte & möchte & would like \\ \hline
\end{tabular}
\end{center}

\begin{warnbox}{Warnung}
\textbf{wollen} im Konjunktiv II ist unhöflich! Verwende stattdessen: \textbf{möchte}
\end{warnbox}

\section{Übungen zum Konjunktiv II}
\label{sec:B2KonjIIUebungen}

\subsection*{Übung 1: Höfliche Form}
\begin{enumerate}
    \item Komm bitte! $\rightarrow$ \underline{\hspace{2cm}}
    \item Hilf mir! $\rightarrow$ \underline{\hspace{2cm}}
    \item Ich will einen Kaffee. $\rightarrow$ \underline{\hspace{2cm}}
    \item Sagen Sie mir die Wahrheit! $\rightarrow$ \underline{\hspace{2cm}}
    \item Geben Sie mir das Buch! $\rightarrow$ \underline{\hspace{2cm}}
\end{enumerate}

\loesung{
\begin{enumerate}
    \item Könntest du bitte kommen?
    \item Würdest du mir helfen?
    \item Ich möchte einen Kaffee.
    \item Würden Sie mir die Wahrheit sagen?
    \item Würden Sie mir das Buch geben?
\end{enumerate}
}

\chapterend

\chapter{C1: Komplexe Satzstrukturen}
\label{cha:C1Komplex}

\section{Parenthesen (Einschübe)}
\label{sec:C1Parenthese}

Parenthesen sind erklärende Einschübe im Satz:

\begin{center}
\begin{tabular}{|l|c|}
\hline
Beispiel & Erklärung \\ \hline
Der Mann - ich kenne ihn seit Jahren - arbeitet bei uns. & Relativsatz als Parenthese \\ \hline
Das Auto, um das es geht, ist quite teuer. & Nomenphrase als Parenthese \\ \hline
\end{tabular}
\end{center}

\section{Attributive Nebensätze}
\label{sec:C1Attributiv}

\begin{center}
\begin{tabular}{|l|c|}
\hline
Nebensatz & Attributsatz \\ \hline
Das ist das Buch, \textbf{das ich gekauft habe}. & Das gekaufte Buch \\ \hline
Die Frau, \textbf{die dort steht}, ist meine Mutter. & Die dort stehende Frau \\ \hline
\end{tabular}
\end{center}

\section{Infinitivgruppen}
\label{sec:C1Infinitiv}

\begin{center}
\begin{tabular}{|l|c|}
\hline
Typ & Beispiel \\ \hline
um zu + Infinitiv & Ich lerne, \textbf{um} eine Prüfung zu bestehen. \\ \hline
ohne zu + Infinitiv & Er ging \textbf{ohne} ein Wort zu sagen. \\ \hline
statt zu + Infinitiv & Sie blieb zu Hause \textbf{statt} zu arbeiten. \\ \hline
\end{tabular}
\end{center}

\section{Modales Passiv}
\label{sec:C1ModalPassiv}

\begin{center}
\begin{tabular}{|l|c|c|}
\hline
Modalität & Aktiv & Passiv \\ \hline
müssen & Er muss das machen. & Das muss gemacht werden. \\ \hline
können & Man kann das machen. & Das kann gemacht werden. \\ \hline
sollen & Die Arbeit soll gemacht werden. & - \\ \hline
\end{tabular}
\end{center}

\chapterend